\section{Spectral truncation of shifted powers}
\label{sec:shifted_density_truncation}

Let $I$ be a finite index set. All matrices in this section act on $\C^I$,
and all norms are Euclidean operator norms. The rank is measured on this
space, before tensoring with an identity on any other space.

% Independent spectral ingredients for OpenAI, Polynomial PEPS approximation
% of gapped square-grid ground states, September 24, 2026, 03-patches.tex,
% lines 484--502, eq:patch-head-tail-bounds, at
% adc7f1241b42e322a6451854ab7e4b4c146bf78a.
% No minimizing tuple, entropy tail, regulator growth, tensor contraction,
% or completion of Proposition 4.1 is asserted by these entries.

\begin{definition}[Closed-threshold spectral projection]
    \label{def:closed_threshold_projection}
    \lean{Matrix.spectralProjectionGE}
    \leanok
    For a Hermitian matrix $A$ and $b\in\R$, let
    $\Pi=\mathbf1_{[b,\infty)}(A)$ be the spectral projection onto the
    eigenvalues at least $b$. An eigenvalue equal to $b$ is included.
\end{definition}

\begin{theorem}[Projection and commutation]
    \label{thm:closed_threshold_projection}
    \lean{Matrix.IsHermitian.spectralProjectionGE_eq_cfc,
        Matrix.IsHermitian.isStarProjection_spectralProjectionGE,
        Matrix.IsHermitian.commute_spectralProjectionGE}
    \leanok
    \uses{def:closed_threshold_projection}
    Let $A$ be Hermitian and $b\in\R$. The projection $\Pi$ is the
    functional calculus of $\mathbf1_{[b,\infty)}$ at $A$, and satisfies
    $\Pi^*=\Pi$, $\Pi^2=\Pi$, and $A\Pi=\Pi A$.
\end{theorem}

\begin{proof}
    \leanok
    The spectrum of $A$ is finite, so its indicator is continuous on that
    spectrum. The functional calculus preserves real-valued functions and
    products. The scalar identities $\mathbf1_{[b,\infty)}^2=
        \mathbf1_{[b,\infty)}$ and
    $t\mathbf{1}_{\lbrack b,\infty)}(t)=\mathbf{1}_{\lbrack b,\infty)}(t)t$
    give the assertions.
\end{proof}

\begin{definition}[Head and tail of a shifted power]
    \label{def:shifted_power_head_tail}
    \lean{Matrix.shiftedPowerHead, Matrix.shiftedPowerTail}
    \leanok
    \uses{def:closed_threshold_projection}
    Let $A\ge0$, $b>0$, and $r\in\R$. Write $F=(A+b\Id)^r$ and define
    \begin{align}
        H & =F\Pi, & T & =F(\Id-\Pi).
        \label{eq:shifted_power_split_definition}
    \end{align}
\end{definition}

\begin{theorem}[Functional calculus of the shifted power]
    \label{thm:shifted_power_functional_calculus}
    \lean{Matrix.PosSemidef.add_smul_one_rpow_eq_cfc,
        Matrix.PosSemidef.commute_add_smul_one_rpow_spectralProjectionGE}
    \leanok
    \uses{def:shifted_power_head_tail}
    If $A\ge0$, $b>0$, and $r\in\R$, then $F=f(A)$ for
    $f(t)=(t+b)^r$, and $F\Pi=\Pi F$.
\end{theorem}

\begin{proof}
    \leanok
    \uses{thm:positive_identity_shift, thm:closed_threshold_projection}
    The matrix $A+b\Id$ is positive definite by
    Theorem~\ref{thm:positive_identity_shift}. Composing the functional
    calculus for $t\mapsto t+b$ and $t\mapsto t^r$ gives $F=f(A)$.
    Both $F$ and $\Pi$ are functions of $A$, so they commute.
\end{proof}

\begin{theorem}[Exact head--tail decomposition]
    \label{thm:shifted_power_split}
    \lean{Matrix.shiftedPowerHead_add_shiftedPowerTail}
    \leanok
    \uses{def:shifted_power_head_tail}
    The matrices in~(\ref{eq:shifted_power_split_definition}) satisfy $H+T=F$.
\end{theorem}

\begin{proof}
    \leanok
    Distributivity gives $F\Pi+F(\Id-\Pi)=F$.
\end{proof}

\begin{theorem}[Scalar descriptions of the head and tail]
    \label{thm:shifted_power_truncated_functions}
    \lean{Matrix.PosSemidef.shiftedPowerHead_eq_cfc,
        Matrix.PosSemidef.shiftedPowerTail_eq_cfc}
    \leanok
    \uses{def:shifted_power_head_tail}
    Let $A\ge0$, $b>0$, and $r\in\R$. Then $H=h(A)$ and $T=t(A)$,
    where
    \begin{align}
        h(\lambda) & =
        \begin{cases}(\lambda+b)^r,&\lambda\ge b,\\0,&\lambda<b,\end{cases}
        \label{eq:shifted_power_head_function} \\
        t(\lambda) & =
        \begin{cases}0,&\lambda\ge b,\\(\lambda+b)^r,&\lambda<b.\end{cases}
        \label{eq:shifted_power_tail_function}
    \end{align}
\end{theorem}

\begin{proof}
    \leanok
    \uses{thm:shifted_power_functional_calculus, thm:closed_threshold_projection}
    Multiplication in the functional calculus gives
    $h(\lambda)=(\lambda+b)^r\mathbf1_{[b,\infty)}(\lambda)$.
    Subtraction gives
    $t(\lambda)=(\lambda+b)^r-h(\lambda)$, yielding the two cases.
\end{proof}

\begin{theorem}[Positivity and orthogonality]
    \label{thm:shifted_power_positive_orthogonal}
    \lean{Matrix.PosSemidef.shiftedPowerHead_posSemidef,
        Matrix.PosSemidef.shiftedPowerTail_posSemidef,
        Matrix.PosSemidef.shiftedPowerHead_mul_shiftedPowerTail,
        Matrix.PosSemidef.shiftedPowerTail_mul_shiftedPowerHead}
    \leanok
    \uses{def:shifted_power_head_tail}
    For $A\ge0$, $b>0$, and every $r\in\R$, one has
    $H\ge0$, $T\ge0$, and $HT=TH=0$.
\end{theorem}

\begin{proof}
    \leanok
    \uses{thm:shifted_power_truncated_functions}
    On the spectrum of $A$, $\lambda+b>0$. Thus both functions in
    (\ref{eq:shifted_power_head_function}) and
    (\ref{eq:shifted_power_tail_function}) are nonnegative, and their
    pointwise product is zero. The functional calculus gives positivity
    and $HT=0$. Taking adjoints gives $TH=0$.
\end{proof}

\begin{theorem}[Inverse shifted powers]
    \label{thm:shifted_power_inverse}
    \lean{Matrix.PosSemidef.add_smul_one_rpow_mul_neg}
    \leanok
    If $A\ge0$, $b>0$, and $r\in\R$, then
    $(A+b\Id)^r(A+b\Id)^{-r}=\Id$.
\end{theorem}

\begin{proof}
    \leanok
    \uses{thm:positive_identity_shift}
    Strict positivity of $A+b\Id$ permits multiplication of its real
    powers. The exponents add to zero.
\end{proof}

\begin{theorem}[Range and rank of the head]
    \label{thm:shifted_power_head_range}
    \lean{Matrix.PosSemidef.range_shiftedPowerHead,
        Matrix.PosSemidef.rank_shiftedPowerHead}
    \leanok
    \uses{def:shifted_power_head_tail}
    For $A\ge0$, $b>0$, and $r\in\R$, the head satisfies
    $\operatorname{ran}H=\operatorname{ran}\Pi$ and
    $\operatorname{rank}H=\operatorname{rank}\Pi$.
\end{theorem}

\begin{proof}
    \leanok
    \uses{thm:shifted_power_functional_calculus, thm:shifted_power_inverse}
    Commutation gives $H=\Pi F$, so
    $\operatorname{ran}H\subseteq\operatorname{ran}\Pi$.
    The inverse identity yields $\Pi=H(A+b\Id)^{-r}$, proving the reverse
    inclusion. Taking dimensions gives equality of ranks.
\end{proof}

\begin{theorem}[Spectral counting and trace]
    \label{thm:closed_threshold_rank_trace}
    \lean{Matrix.PosSemidef.mul_rank_spectralProjectionGE_le_trace}
    \leanok
    \uses{def:closed_threshold_projection}
    For every $A\ge0$ and every real $b$,
    $b\operatorname{rank}\Pi\le\operatorname{Re}\tr A$.
\end{theorem}

\begin{proof}
    \leanok
    \uses{thm:closed_threshold_projection}
    The rank of an orthogonal projection equals its trace. If
    $\lambda_i$ are the eigenvalues of $A$, then
    \begin{align}
        b\operatorname{rank}\Pi
        =\sum_{i\in I}b\mathbf1_{[b,\infty)}(\lambda_i)
        \le\sum_{i\in I}\lambda_i=\operatorname{Re}\tr A.
    \end{align}
    Each scalar inequality follows from $b\le\lambda_i$ on the selected
    spectrum and $0\le\lambda_i$ on its complement.
\end{proof}

\begin{theorem}[Head rank for a density matrix]
    \label{thm:shifted_density_head_rank}
    \lean{Matrix.PosSemidef.rank_shiftedPowerHead_le}
    \leanok
    \uses{def:shifted_power_head_tail}
    If $A\ge0$, $\tr A=1$, $b>0$, and $r\in\R$, then
    $\operatorname{rank}H\le b^{-1}$.
\end{theorem}

\begin{proof}
    \leanok
    \uses{thm:shifted_power_head_range, thm:closed_threshold_rank_trace}
    Substitute $\operatorname{rank}H=\operatorname{rank}\Pi$ into the
    trace inequality and divide by $b>0$.
\end{proof}

\begin{theorem}[Tail norm for nonnegative powers]
    \label{thm:shifted_power_tail_norm}
    \lean{Matrix.PosSemidef.l2_opNorm_shiftedPowerTail_le}
    \leanok
    \uses{def:shifted_power_head_tail}
    If $A\ge0$, $b>0$, and $r\ge0$, then $\|T\|\le(2b)^r$.
    No trace normalization is required.
\end{theorem}

\begin{proof}
    \leanok
    \uses{thm:shifted_power_truncated_functions}
    The tail function vanishes for $\lambda\ge b$. For the remaining
    spectral values, $0\le\lambda<b$, so $0<\lambda+b<2b$ and
    $(\lambda+b)^r\le(2b)^r$ by monotonicity for $r\ge0$.
    The norm estimate for functional calculus gives the result.
\end{proof}

For $A\ge0$, $b>0$, and $r\in\R$, consider the orthogonal decomposition
$\C^I=\operatorname{ran}\Pi\oplus\ker\Pi$. Let $F_{\ge b}$ and $F_{<b}$
be the restrictions of $F$ to the two summands. The split has the form
\begin{align}
    H+T
    =\begin{pmatrix}F_{\ge b}&0\\0&0\end{pmatrix}
    +\begin{pmatrix}0&0\\0&F_{<b}\end{pmatrix}
    =\begin{pmatrix}F_{\ge b}&0\\0&F_{<b}\end{pmatrix}=F.
\end{align}
The restriction $F_{\ge b}$ is invertible on $\operatorname{ran}\Pi$,
which explains why truncating the power preserves the projection's rank.

If $I$ is empty, the matrices, ranges, and ranks are zero, and the stated
identities and estimates remain valid. A trace-one matrix cannot occur in
that case. The exponent $r=0$ gives $F=\Id$, $H=\Pi$, and $T=\Id-\Pi$.
